\section{Nested Datatypes: Bush}
\label{sec:bush}

One motivation for nested datatypes is the representation of function types as terminal coalgebras, that is, as coinductive types. A well-known instance is \AgdaDatatype{ℕ}\;\AgdaSymbol{→}\;\AgdaBound{X}\;\AgdaFunction{≃}\;\AgdaRecord{Stream}\;\AgdaBound{X}, where \AgdaRecord{Stream}\;\AgdaBound{X} is defined coinductively by \AgdaBound{X}\;\AgdaFunction{×}\;\AgdaRecord{Stream}\;\AgdaBound{X}. Now consider binary trees carrying no data at either leaves or nodes:

\begin{code}%
\>[0]\AgdaKeyword{data}\AgdaSpace{}%
\AgdaDatatype{BTree}\AgdaSpace{}%
\AgdaSymbol{:}\AgdaSpace{}%
\AgdaPrimitive{Type}\AgdaSpace{}%
\AgdaKeyword{where}\<%
\\
\>[0][@{}l@{\AgdaIndent{0}}]%
\>[2]\AgdaInductiveConstructor{leaf}\AgdaSpace{}%
\AgdaSymbol{:}\AgdaSpace{}%
\AgdaDatatype{BTree}\<%
\\
%
\>[2]\AgdaInductiveConstructor{node}\AgdaSpace{}%
\AgdaSymbol{:}\AgdaSpace{}%
\AgdaDatatype{BTree}\AgdaSpace{}%
\AgdaSymbol{→}\AgdaSpace{}%
\AgdaDatatype{BTree}\AgdaSpace{}%
\AgdaSymbol{→}\AgdaSpace{}%
\AgdaDatatype{BTree}\<%
\end{code}


The corresponding terminal coalgebra of \AgdaDatatype{BTree}\;\AgdaSymbol{→}\;\AgdaBound{X} is given by the following type:

\begin{code}
\>[0]\AgdaKeyword{record}\AgdaSpace{}%
\AgdaRecord{Bush∞}\AgdaSpace{}%
\AgdaSymbol{(}\AgdaBound{X}\AgdaSpace{}%
\AgdaSymbol{:}\AgdaSpace{}%
\AgdaPrimitive{Type}\AgdaSymbol{)}\AgdaSpace{}%
\AgdaSymbol{:}\AgdaSpace{}%
\AgdaPrimitive{Type}\AgdaSpace{}%
\AgdaKeyword{where}\<%
\\
\>[0][@{}l@{\AgdaIndent{0}}]%
\>[2]\AgdaKeyword{coinductive}\<%
\\
%
\>[2]\AgdaKeyword{field}\<%
\\
\>[2][@{}l@{\AgdaIndent{0}}]%
\>[4]\AgdaField{head}\AgdaSpace{}%
\AgdaSymbol{:}\AgdaSpace{}%
\AgdaBound{X}\<%
\\
%
\>[4]\AgdaField{tail}\AgdaSpace{}%
\AgdaSymbol{:}\AgdaSpace{}%
\AgdaRecord{Bush∞}\AgdaSpace{}%
\AgdaSymbol{(}\AgdaRecord{Bush∞}\AgdaSpace{}%
\AgdaBound{X}\AgdaSymbol{)}\<%
\end{code}


The destructor \AgdaField{tail} involves a double application of \AgdaRecord{Bush∞}. In fact, nesting of arbitrary finite depth is required: one can show that the function type out of $n$-ary branching trees \AgdaDatatype{NTree}\;\AgdaSymbol{→}\;\AgdaBound{X} corresponds to a terminal coalgebra \AgdaRecord{Bushⁿ∞} whose destructor \AgdaField{tail} applies \AgdaRecord{Bushⁿ∞} $n$ times. For simplicity and intuition, however, we first consider an inhabited, inductive variant:

\input{code/bush}

This type closely resembles \AgdaDatatype{List}, except for the nested application in \AgdaInductiveConstructor{\_∷\_}. Whereas the signature of \AgdaDatatype{List}\;\AgdaBound{A} can be presented as \AgdaBound{X}\;\AgdaSymbol{↦}\;\AgdaRecord{⊤}\;\AgdaDatatype{⊎}\;\AgdaBound{A}\;\AgdaFunction{×}\;\AgdaBound{X}, the signature of \AgdaDatatype{Bush}\;\AgdaBound{A} cannot be brought into the same form \AgdaBound{X}\;\AgdaSymbol{↦}\;\AgdaRecord{⊤}\;\AgdaDatatype{⊎}\;\AgdaBound{A}\;\AgdaFunction{×}\;\AgdaBound{?} since there is no type that can be placed in the hole \AgdaBound{?} to express the double application. Our solution is to lift the setting to second order: instead of algebras in the category of types, we consider algebras in the category of containers. At the level of types, a second-order signature for \AgdaDatatype{Bush} is given by:

\begin{code}%
\>[0]\AgdaFunction{SigBush}\AgdaSpace{}%
\AgdaSymbol{:}\AgdaSpace{}%
\AgdaSymbol{(}\AgdaPrimitive{Type}\AgdaSpace{}%
\AgdaSymbol{→}\AgdaSpace{}%
\AgdaPrimitive{Type}\AgdaSymbol{)}\AgdaSpace{}%
\AgdaSymbol{→}\AgdaSpace{}%
\AgdaPrimitive{Type}\AgdaSpace{}%
\AgdaSymbol{→}\AgdaSpace{}%
\AgdaPrimitive{Type}\<%
\\
\>[0]\AgdaFunction{SigBush}\AgdaSpace{}%
\AgdaBound{F}\AgdaSpace{}%
\AgdaBound{X}\AgdaSpace{}%
\AgdaSymbol{=}\AgdaSpace{}%
\AgdaRecord{⊤}\AgdaSpace{}%
\AgdaOperator{\AgdaDatatype{⊎}}\AgdaSpace{}%
\AgdaBound{X}\AgdaSpace{}%
\AgdaOperator{\AgdaFunction{×}}\AgdaSpace{}%
\AgdaBound{F}\AgdaSpace{}%
\AgdaSymbol{(}\AgdaBound{F}\AgdaSpace{}%
\AgdaBound{X}\AgdaSymbol{)}\<%
\end{code}


This signature abstracts over a function variable \AgdaBound{F}. If we restrict \AgdaBound{F} to be strictly positive, that is, representable by a container, the signature can be expressed as an endofunctor on containers:

\begin{code}%
\>[0]\AgdaFunction{SigBushᶜ}\AgdaSpace{}%
\AgdaSymbol{:}\AgdaSpace{}%
\AgdaRecord{Cont}\AgdaSpace{}%
\AgdaSymbol{→}\AgdaSpace{}%
\AgdaRecord{Cont}\<%
\\
\>[0]\AgdaFunction{SigBushᶜ}\AgdaSpace{}%
\AgdaBound{SP}\AgdaSpace{}%
\AgdaSymbol{=}\AgdaSpace{}%
\AgdaFunction{⊤ᶜ}\AgdaSpace{}%
\AgdaOperator{\AgdaFunction{⊎ᶜ}}\AgdaSpace{}%
\AgdaSymbol{(}\AgdaFunction{1ᶜ}\AgdaSpace{}%
\AgdaOperator{\AgdaFunction{×ᶜ}}\AgdaSpace{}%
\AgdaSymbol{(}\AgdaBound{SP}\AgdaSpace{}%
\AgdaOperator{\AgdaFunction{⊗ᶜ}}\AgdaSpace{}%
\AgdaBound{SP}\AgdaSymbol{))}\<%
\end{code}


We now construct a container representation of \AgdaDatatype{Bush}. Suppose that \AgdaFunction{Bushᶜ}\;\AgdaSymbol{=}\;\AgdaRecord{S}\;\AgdaInductiveConstructor{◃}\;\AgdaRecord{P} is the initial algebra of \AgdaFunction{SigBushᶜ}. The construction is guided by the following calculation:

\[
\begin{array}{cl}
  & \AgdaRecord{S}\;\AgdaInductiveConstructor{◃}\;\AgdaRecord{P} \\
  & \quad \{\ \text{by assumption} \} \\
≃ & \AgdaFunction{SigBushᶜ}\;(\AgdaRecord{S}\;\AgdaInductiveConstructor{◃}\;\AgdaRecord{P}) \\
  & \quad \{\ \text{by definition of \AgdaFunction{SigBushᶜ}}\ \} \\
≃ & \AgdaFunction{⊤ᶜ}\;\AgdaFunction{⊎ᶜ}\;\AgdaFunction{1ᶜ}\;\AgdaFunction{×ᶜ}\;(\AgdaRecord{S}\;\AgdaInductiveConstructor{◃}\;\AgdaRecord{P})\;\AgdaOperator{\AgdaFunction{⊗ᶜ}}\;(\AgdaRecord{S}\;\AgdaInductiveConstructor{◃}\;\AgdaRecord{P}) \\
  & \quad \{\ \text{by definition of \AgdaOperator{\AgdaFunction{⊗ᶜ}}}\ \} \\
≃ & \AgdaFunction{⊤ᶜ}\;\AgdaFunction{⊎ᶜ}\;\AgdaFunction{1ᶜ}\;\AgdaFunction{×ᶜ}\;(\AgdaFunction{Σ}\,[\,\AgdaBound{s}\in\AgdaRecord{S}\,]\;(\AgdaRecord{P}\;\AgdaBound{s}\to\AgdaRecord{S})
    \;\AgdaInductiveConstructor{◃}\;
\AgdaSymbol{λ\{}\;(\AgdaBound{s}\;\AgdaOperator{\AgdaInductiveConstructor{,}}\;\AgdaBound{f})\to\AgdaFunction{Σ}\,[\,\AgdaBound{p}\in\AgdaRecord{P}\;\AgdaBound{s}\,]\;\AgdaRecord{P}\,(\AgdaBound{f}\;\AgdaBound{p})\AgdaSymbol{\}}) \\
  & \quad \{\ \text{by definition of \AgdaFunction{1ᶜ} and \AgdaFunction{×ᶜ}}\ \} \\
≃ & \AgdaFunction{⊤ᶜ}\;\AgdaFunction{⊎ᶜ}\;(\AgdaFunction{Σ}\,[\,\AgdaBound{s}\in\AgdaRecord{S}\,]\;(\AgdaRecord{P}\;\AgdaBound{s}\to\AgdaRecord{S})
    \;\AgdaInductiveConstructor{◃}\;
\AgdaSymbol{λ\{}\;(\AgdaBound{s}\;\AgdaOperator{\AgdaInductiveConstructor{,}}\;\AgdaBound{f})\to\AgdaDatatype{⊤}\;\AgdaDatatype{⊎}\;\AgdaFunction{Σ}\,[\,\AgdaBound{p}\in\AgdaRecord{P}\;\AgdaBound{s}\,]\;\AgdaRecord{P}\,(\AgdaBound{f}\;\AgdaBound{p})\AgdaSymbol{\}}) \\
  & \quad \{\ \text{by definition of \AgdaFunction{⊤ᶜ} and \AgdaFunction{⊎ᶜ}}\ \} \\
≃ & \AgdaDatatype{⊤}\;\AgdaDatatype{⊎}\;\AgdaFunction{Σ}\,[\,\AgdaBound{s}\in\AgdaRecord{S}\,]\;(\AgdaRecord{P}\;\AgdaBound{s}\to\AgdaRecord{S})
    \;\AgdaInductiveConstructor{◃}\;
\AgdaSymbol{λ\{}\;(\AgdaInductiveConstructor{inj₁}\;\AgdaInductiveConstructor{tt})\to\AgdaDatatype{⊥}\;\AgdaSymbol{;}\\
  & \hphantom{\AgdaDatatype{⊤}\;\AgdaDatatype{⊎}\;\AgdaFunction{Σ}\,[\,\AgdaBound{s}\in\AgdaRecord{S}\,]\;(\AgdaRecord{P}\;\AgdaBound{s}\to\AgdaRecord{S})\;\AgdaInductiveConstructor{◃}\;\AgdaSymbol{λ\{}}\;
(\AgdaInductiveConstructor{inj₂}\;(\AgdaBound{s}\;\AgdaOperator{\AgdaInductiveConstructor{,}}\;\AgdaBound{f}))\to\AgdaDatatype{⊤}\;\AgdaDatatype{⊎}\;\AgdaFunction{Σ}\,[\,\AgdaBound{p}\in\AgdaRecord{P}\;\AgdaBound{s}\,]\;\AgdaRecord{P}\,(\AgdaBound{f}\;\AgdaBound{p})\AgdaSymbol{\}}
\end{array}
\]


Reading off the solution of this equation, we define the shapes and positions mutually by induction-induction:

\begin{code}%
\>[2][@{}l@{\AgdaIndent{1}}]%
\>[4]\AgdaKeyword{record}\AgdaSpace{}%
\AgdaRecord{S}\AgdaSpace{}%
\AgdaSymbol{:}\AgdaSpace{}%
\AgdaPrimitive{Type}\AgdaSpace{}%
\AgdaKeyword{where}\<%
\\
\>[4][@{}l@{\AgdaIndent{0}}]%
\>[6]\AgdaKeyword{inductive}\<%
\\
%
\>[6]\AgdaKeyword{field}\<%
\\
\>[6][@{}l@{\AgdaIndent{0}}]%
\>[8]\AgdaField{desS}\AgdaSpace{}%
\AgdaSymbol{:}\AgdaSpace{}%
\AgdaRecord{⊤}\AgdaSpace{}%
\AgdaOperator{\AgdaDatatype{⊎}}\AgdaSpace{}%
\AgdaFunction{Σ[}\AgdaSpace{}%
\AgdaBound{s}\AgdaSpace{}%
\AgdaFunction{∈}\AgdaSpace{}%
\AgdaRecord{S}\AgdaSpace{}%
\AgdaFunction{]}\AgdaSpace{}%
\AgdaSymbol{(}\AgdaRecord{P}\AgdaSpace{}%
\AgdaBound{s}\AgdaSpace{}%
\AgdaSymbol{→}\AgdaSpace{}%
\AgdaRecord{S}\AgdaSymbol{)}\<%
\\
%
\\[\AgdaEmptyExtraSkip]%
%
\>[4]\AgdaKeyword{record}\AgdaSpace{}%
\AgdaRecord{P}\AgdaSpace{}%
\AgdaSymbol{(}\AgdaBound{s}\AgdaSpace{}%
\AgdaSymbol{:}\AgdaSpace{}%
\AgdaRecord{S}\AgdaSymbol{)}\AgdaSpace{}%
\AgdaSymbol{:}\AgdaSpace{}%
\AgdaPrimitive{Type}\AgdaSpace{}%
\AgdaKeyword{where}\<%
\\
\>[4][@{}l@{\AgdaIndent{0}}]%
\>[6]\AgdaKeyword{inductive}\<%
\\
%
\>[6]\AgdaKeyword{field}\<%
\\
\>[6][@{}l@{\AgdaIndent{0}}]%
\>[8]\AgdaField{desP}\AgdaSpace{}%
\AgdaSymbol{:}\AgdaSpace{}%
\AgdaFunction{case}\AgdaSpace{}%
\AgdaSymbol{(λ\{}\AgdaSpace{}%
\AgdaInductiveConstructor{tt}\AgdaSpace{}%
\AgdaSymbol{→}\AgdaSpace{}%
\AgdaDatatype{⊥}\AgdaSpace{}%
\AgdaSymbol{\})}\AgdaSpace{}%
\AgdaSymbol{(λ\{}\AgdaSpace{}%
\AgdaSymbol{(}\AgdaBound{s}\AgdaSpace{}%
\AgdaOperator{\AgdaInductiveConstructor{,}}\AgdaSpace{}%
\AgdaBound{f}\AgdaSymbol{)}\AgdaSpace{}%
\AgdaSymbol{→}\AgdaSpace{}%
\AgdaRecord{⊤}\AgdaSpace{}%
\AgdaOperator{\AgdaDatatype{⊎}}\AgdaSpace{}%
\AgdaFunction{Σ[}\AgdaSpace{}%
\AgdaBound{p}\AgdaSpace{}%
\AgdaFunction{∈}\AgdaSpace{}%
\AgdaRecord{P}\AgdaSpace{}%
\AgdaBound{s}\AgdaSpace{}%
\AgdaFunction{]}\AgdaSpace{}%
\AgdaRecord{P}\AgdaSpace{}%
\AgdaSymbol{(}\AgdaBound{f}\AgdaSpace{}%
\AgdaBound{p}\AgdaSymbol{)}\AgdaSpace{}%
\AgdaSymbol{\})}\AgdaSpace{}%
\AgdaSymbol{(}\AgdaField{S.desS}\AgdaSpace{}%
\AgdaBound{s}\AgdaSymbol{)}\<%
\\
%
\\[\AgdaEmptyExtraSkip]%
%
\>[2]\AgdaFunction{Bushᶜ}\AgdaSpace{}%
\AgdaSymbol{:}\AgdaSpace{}%
\AgdaRecord{Cont}\<%
\\
%
\>[2]\AgdaFunction{Bushᶜ}\AgdaSpace{}%
\AgdaSymbol{=}\AgdaSpace{}%
\AgdaRecord{S}\AgdaSpace{}%
\AgdaOperator{\AgdaInductiveConstructor{◃}}\AgdaSpace{}%
\AgdaRecord{P}\<%
\end{code}


The shape \AgdaRecord{S} is inductive because we are constructing an initial algebra; to obtain the terminal coalgebra, it suffices to make it coinductive instead. The position \AgdaRecord{P}, however, should be always inductive, since positions must be finite: you only provide finite terms to a constructor or a field. We prove that \AgdaFunction{Bushᶜ} is indeed the initial algebra of \AgdaFunction{SigBushᶜ}; the proof can be found in our formalization. Finally, the raw \AgdaDatatype{Bush} type, together with its functoriality, can be derived by calling the type interpretation \AgdaFunction{⟦\_⟧} on \AgdaFunction{Bushᶜ}.
